\documentclass[runningheads]{llncs}
\usepackage{paper-eccv}	% see this file for paper ID and other options

\begin{document}

\title{Paper Title}

\maketitle

\begin{abstract}
	Your abstract.
	\keywords{First keyword \and Second keyword \and Third keyword}
\end{abstract}

\begin{figure}[t]
	\centering
	\includegraphics[width=\textwidth]{teaser.pdf}
	\caption{Your teaser. Replace \texttt{figures/teaser.pdf} with a figure that shows the result of the paper, and write a caption that stands on its own.}
	\label{fig:teaser}%
\end{figure}

\section{Introduction}
\label{sec:introduction}

Use this document as a template for ECCV papers. Representative citation of a cool paper \citep{Miller:VOS:2024}.

\section{Notation and results}
\label{sec:results}

Notation macros belong in \texttt{paper-eccv.sty}, so that changing a symbol is one edit. Equations are typeset with the packages and macros \texttt{mathdefs} provides,
%
\begin{equation}
	\label{equ:example}
	\mathboxnt{boxbg3}{\E{f \paren{\bm{x}}}} = \mathbox{boxtitle1}{boxbg1}{integral}{\int_{\R^3} f \paren{\bm{x}} \, p \paren{\bm{x}} \ud \bm{x}},
\end{equation}
%
and \texttt{boxdefs} provides boxes to highlight their terms, in colors from the \texttt{colordefs} palette, as well as environments for the statements you want to stand out.

\begin{prp}{An example proposition}{example}
	State the claim here. The environment numbers propositions, definitions, and lemmas in a single sequence, and labels this one \texttt{pro:example}.
\end{prp}

\Cref{pro:example} follows from \cref{equ:example}, which is the sort of cross-reference \texttt{cleveref} writes out for you.

\begin{algorithm}[t]
	\caption{\Proc{EstimateIntegral}$(f, p, N)$}
	\label{alg:estimate}
	\begin{algorithmic}[1]
		\Require \textit{Integrand $f$, density $p$, sample count $N$.}
		\Ensure \textit{Estimate $\widehat{F}$ of the integral in \cref{equ:example}.}
		\State $\widehat{F} \gets 0$
		\FORDO[$n \gets 1, \dots, N$]
			\COMMENT[Sample a point from the density.]\ENDCOMMENT
			\State $\bm{x}_n \gets \Proc{Sample}(p)$ \label{alg:estimate:sample}
			\COMMENT[Evaluate the integrand at that point.]\ENDCOMMENT
			\State $\widehat{F} \gets \widehat{F} + f \sparen{\bm{x}_n}$ \label{alg:estimate:evaluate}
		\ENDFORDO
		\RETURN[$\widehat{F} / N$]\ENDRETURN
	\end{algorithmic}
\end{algorithm}

Algorithms are typeset with the packages and macros \texttt{algodefs} provides. \Cref{alg:estimate} estimates the integral in \cref{equ:example} by sampling the density and evaluating the integrand at each sample (\cref{alg:estimate:sample,alg:estimate:evaluate}).

\bibliographystyle{splncs04nat}
\bibliography{paper-eccv}

\end{document}
